\documentclass[11pt]{article}
\usepackage{amsmath,amssymb,amsthm}
\newtheorem{theorem}{Theorem}
\newtheorem{lemma}[theorem]{Lemma}
\newtheorem{remark}[theorem]{Remark}
\begin{document}

\noindent
Let \(T,T'\in B(H)\) be isometries with \(TT'=T'T\).
Write \(D_{T^*}=(I-TT^*)^{1/2}\) and \(D_{{T'}^*}=(I-T'T'^*)^{1/2}\).
Since \(T\) is an isometry one has \(D_T=0\) and \(D_{T^*}=I-TT^*=P_{\ker T^*}\)
(and likewise for \(T'\)).

%% ============================================================
\section*{1. Single isometry: Halmos block matrix}

On \(K_1=H\oplus H\) define
\begin{equation}
\label{eq:Halmos}
\boxed{
U
=
\begin{pmatrix}
T & D_{T^*} \\
0 & -T^*
\end{pmatrix}.
}
\end{equation}

\begin{lemma}
\(U\) is unitary and \(U\big|_{H\oplus\{0\}}=T\). In particular \(U\) is a unitary
\emph{extension} of \(T\), and \(U^n\big|_H=T^n\) for all \(n\ge 0\).
\end{lemma}

\begin{proof}
The identity \(TD_T=D_{T^*}T\) collapses to \(D_{T^*}T=0\). Hence
\begin{align*}
U^*U
&=
\begin{pmatrix}
T^* & 0 \\
D_{T^*} & -T
\end{pmatrix}
\begin{pmatrix}
T & D_{T^*} \\
0 & -T^*
\end{pmatrix}
=
\begin{pmatrix}
T^*T & T^*D_{T^*} \\
D_{T^*}T & D_{T^*}^2+TT^*
\end{pmatrix}
=
\begin{pmatrix}
I & 0 \\
0 & I
\end{pmatrix},
\end{align*}
so \(U\) is unitary. Also \(U(h,0)=(Th,0)\).
\end{proof}

%% ============================================================
\section*{2. Doubly commuting pair: both Halmos matrices on \(H\oplus H\)}

Assume in addition \(T^*T'=T'T^*\) (equivalently \(TD_{{T'}^*}=D_{{T'}^*}T\) and
\(T'D_{T^*}=D_{T^*}T'\)).

On the \emph{same} space \(H\oplus H\) set
\begin{equation}
\label{eq:both-Halmos}
\boxed{
U
=
\begin{pmatrix}
T & D_{T^*} \\
0 & -T^*
\end{pmatrix},
\qquad
U'
=
\begin{pmatrix}
T' & D_{{T'}^*} \\
0 & -{T'}^*
\end{pmatrix}.
}
\end{equation}

\begin{theorem}
If \(T,T'\) doubly commute, then \(U\) and \(U'\) are commuting unitaries on
\(H\oplus H\) extending \(T\) and \(T'\) respectively.
\end{theorem}

\begin{proof}
Each factor is unitary by Section~1. Compute
\begin{align*}
UU'
&=
\begin{pmatrix}
TT' & T\,D_{{T'}^*}-D_{T^*}{T'}^* \\
0 & T^*{T'}^*
\end{pmatrix},
&
U'U
&=
\begin{pmatrix}
T'T & T'D_{T^*}-D_{{T'}^*}T^* \\
0 & {T'}^*T^*
\end{pmatrix}.
\end{align*}
Double commutativity gives \(TT'=T'T\), \(T^*{T'}^*={T'}^*T^*\), and
\[
T\,D_{{T'}^*}-D_{T^*}{T'}^*
=
T'D_{T^*}-D_{{T'}^*}T^*,
\]
so \(UU'=U'U\).
\end{proof}

%% ============================================================
\section*{3. General commuting pair: Halmos, then extend the commutant}

Now assume only \(TT'=T'T\) (not necessarily double).

\subsection*{Step A --- Halmos extension of \(T\)}
On \(K_1=H\oplus H\),
\[
U
=
\begin{pmatrix}
T & D_{T^*} \\
0 & -T^*
\end{pmatrix}.
\]

\subsection*{Step B --- isometric extension of \(T'\) commuting with \(U\)}
Define \(\widetilde{T}'\) on the dense manifold
\(\operatorname{span}\{U^n(h,0):n\in\mathbb{Z},\,h\in H\}\) by
\begin{equation}
\label{eq:inductive}
\widetilde{T}'\bigl(U^n(h,0)\bigr)=U^n(T'h,0).
\end{equation}
The same inner-product computation as in the classical inductive proof shows that
\(\widetilde{T}'\) is a well-defined isometry on all of \(K_1\) (for the Halmos \(U\) of an
isometry one has \(\bigvee_{n\in\mathbb{Z}}U^n(H\oplus 0)=H\oplus H\)), that
\[
\widetilde{T}'\big|_{H\oplus 0}=T',
\qquad
\widetilde{T}'U=U\widetilde{T}',
\]
and that \(\widetilde{T}'\) is unitary whenever \(T'\) is.

\medskip
\noindent
\textbf{Explicit \(2\times 2\) block form of \(\widetilde{T}'\).}
Write vectors as \((x,y)\in H\oplus H\). One checks from~\eqref{eq:inductive} and
\(U^*(h,0)=(T^*h,\,D_{T^*}h)\) that
\[
\widetilde{T}'
=
\begin{pmatrix}
T' & X \\
0 & Y
\end{pmatrix},
\]
where the operators \(X,Y\in B(H)\) are uniquely determined by: for all \(d\in\ker T^*=\operatorname{Ran}D_{T^*}\),
\begin{equation}
\label{eq:XY-on-defect}
Xd=0,
\qquad
Yd=D_{T^*}T'd,
\end{equation}
together with the commutation identities forced by \(\widetilde{T}'U=U\widetilde{T}'\),
\begin{equation}
\label{eq:XY-comm}
YT^*=T^*Y,
\qquad
T'D_{T^*}-XT^*=TX+D_{T^*}Y,
\end{equation}
and the isometry conditions
\begin{equation}
\label{eq:XY-isom}
T'^*X=0,
\qquad
X^*X+Y^*Y=I.
\end{equation}
(When \(T,T'\) doubly commute one may take \(X=D_{{T'}^*}\) and \(Y=-{T'}^*\), recovering~\eqref{eq:both-Halmos};
in general \(X,Y\) are the unique solution of~\eqref{eq:XY-on-defect}--\eqref{eq:XY-isom} obtained from~\eqref{eq:inductive}.)

\subsection*{Step C --- Halmos unitarization of \(\widetilde{T}'\) (if needed)}
If \(\widetilde{T}'\) is already unitary, stop: \((U,\widetilde{T}')\) is the desired pair on \(H\oplus H\).

If not, apply Halmos to \(\widetilde{T}'\) on \(K=K_1\oplus K_1=(H\oplus H)\oplus(H\oplus H)\):
\begin{equation}
\label{eq:Uprime-Halmos}
\boxed{
U'
=
\begin{pmatrix}
\widetilde{T}' & D_{\widetilde{T}'{}^*} \\
0 & -{\widetilde{T}'}^*
\end{pmatrix}.
}
\end{equation}
Extend \(U\) to a unitary \(\widetilde{U}\) on \(K\) by the same inductive rule as~\eqref{eq:inductive}
(with \((\widetilde{T}',U)\) in place of \((T,T')\)). Since \(U\) is already unitary, the resulting
\(\widetilde{U}\) is unitary and
\[
\widetilde{U}\,U'=U'\,\widetilde{U}.
\]
Thus \((\widetilde{U},U')\) is a commuting unitary extension of \((T,T')\) on
\[
K=H\oplus H\oplus H\oplus H.
\]

\begin{remark}[block \(4\times 4\) shape]
Writing \(K\) as four copies of \(H\), the operators have the schematic form
\[
U'
=
\begin{pmatrix}
T' & X & * & * \\
0 & Y & * & * \\
0 & 0 & -Y^* & * \\
0 & 0 & -X^* & -{T'}^*
\end{pmatrix},
\qquad
\widetilde{U}
=
\begin{pmatrix}
T & D_{T^*} & * & * \\
0 & -T^* & * & * \\
0 & 0 & * & * \\
0 & 0 & * & *
\end{pmatrix},
\]
where \(X,Y\) are as in Step~B and the remaining blocks are those of
\(D_{\widetilde{T}'{}^*}\) and the inductive extension of \(U\).
\end{remark}

%% ============================================================
\section*{4. Summary}

\begin{center}
\renewcommand{\arraystretch}{1.4}
\begin{tabular}{ll}
\hline
single isometry &
\(\displaystyle U=\begin{pmatrix}T&D_{T^*}\\ 0&-T^*\end{pmatrix}\) on \(H\oplus H\) \\
doubly commuting pair &
both Halmos matrices~\eqref{eq:both-Halmos} on \(H\oplus H\) \\
general commuting pair &
Halmos \(U\) on \(H\oplus H\), then
\(\widetilde{T}'=\begin{pmatrix}T'&X\\ 0&Y\end{pmatrix}\), \\
& then Halmos \(U'\) of \(\widetilde{T}'\) on \((H\oplus H)^{\oplus 2}\) \\
\hline
\end{tabular}
\end{center}

\end{document}
